\SourceLeaf{95}
\Article{120}

Colliguntur facile ex praecc. theoremata haec (vid. artt. 102, 103, 105).

I. $-3$ \emph{est residuum omnium numerorum, qui neque per $8$, neque per $9$, neque per ullum numerum primum formae $6n+5$ dividi possunt, non-residuum autem omnium reliquorum.}

II. $+3$ \emph{est residuum omnium numerorum, qui neque per $4$, neque per $9$, neque per ullum primum formae $12n+5$ vel $12n+7$ dividi possunt, omnium reliquorum non-residuum.}

Teneatur imprimis casus particularis hic:

$-3$ est \emph{residuum} omnium numerorum primorum formae $3n+1$, seu quod idem est \emph{omnium, qui ipsius $3$ sunt residua, non-residuum} vero omnium numerorum primorum formae $6n+5$, seu, excluso numero $2$, omnium formae $3n+2$, \emph{i. e. omnium qui ipsius $3$ sunt non-residua.} Facile vero perspicitur omnes reliquos casus ex hoc sponte sequi.

Propositiones ad residua $+3$ et $-3$ pertinentes iam Fermatio notae fuerunt, \emph{Opera Wallisii T.} II. \emph{p.} 857. At ill. Euler primus demonstrationes tradidit, \emph{Comm. nov. Petr. T.} VIII. \emph{p.} 105 sqq. Eo magis est mirandum, demonstrationes propositionum ad residua $+2$ et $-2$ pertinentium, prorsus similibus artificiis innixas, semper ipsius sagacitatem fugisse. Vid. etiam comment. ill. La Grange, \emph{Nouv. M\'em. de l'Ac. de Berlin}, 1775 \emph{p.} 352.

\begin{center}\emph{Residua $+5$ et $-5$.}\end{center}

\Article{121}

Per inductionem deprehenditur, $+5$ nullius numeri imparis formae $5n+2$ vel $5n+3$ residuum esse, \emph{i. e.} nullius numeri imparis qui ipsius $5$ non-residuum sit. Hanc vero regulam nullam exceptionem pati, ita demonstratur. Sit numerus minimus, si quis datur, ab hac regula excipiendus $=t$, qui itaque numeri $5$ est non-residuum, $5$ autem ipsius $t$ residuum. Sit $aa=5+tu$, ita ut $a$ sit par ipsoque $t$ minor. Erit igitur $u$ impar ipsoque $t$ minor, $+5$ autem ipsius $u$ residuum erit. Quodsi iam $a$ per $5$ non est divisibilis, etiam $u$ non erit; manifesto autem $tu$ ipsius $5$ est residuum, quare quum $t$ ipsius $5$ sit non-residuum, etiam $u$ non-residuum erit; \emph{i. e.} datur non-residuum impar numeri $5$, cuius residuum est $+5$, ipso $t$ minus, contra hyp. Si vero $a$ per $5$ est divisi\SourceLeaf{96}bilis, ponatur $a=5b$, atque $u=5v$, unde $tv\equiv -1\equiv 4\,(\operatorname{mod.}5)$, \emph{i. e.} $tv$ erit residuum numeri 5. In reliquis demonstratio perinde procedit ut in casu priori.

\Article{122}

Omnium igitur numerorum primorum, qui simul sunt ipsius 5 non-residua simulque formae $4n+1$, \emph{i. e.} omnium numerorum primorum formae $20n+13$ vel $20n+17$, tum $+5$ quam $-5$ non-residua erunt; omnium autem numerorum primorum formae $20n+3$ vel $20n+7$, non-residuum erit $+5$, $-5$ residuum.

Potest vero prorsus simili modo demonstrari, $-5$ esse non-residuum omnium numerorum primorum formarum $20n+11$, $20n+13$, $20n+17$, $20n+19$, facileque perspicitur hinc sequi, $+5$ esse residuum omnium numerorum primorum formae $20n+11$ vel $20n+19$, non-residuum autem omnium formae $20n+13$ vel $20n+17$. Et quoniam quivis numerus primus, praeter 2 et 5 (quorum residuum $\pm 5$), in aliqua harum formarum continetur $20n+1,3,7,9,11,13,17,19$, patet, de omnibus iam iudicium ferri posse, exceptis iis qui sint formae $20n+1$ vel formae $20n+9$.

\Article{123}

Ex inductione facile deprehenditur, $+5$ et $-5$ esse residua omnium numerorum primorum formae $20n+1$ vel $20n+9$. Quodsi hoc generaliter verum est, lex elegans habebitur, \emph{$+5$ esse residuum omnium numerorum primorum qui ipsius 5 sint residua} (hi enim in alterutra formarum $5n+1$ vel $5n+4$ sive in aliqua harum, $20n+1,9,11,19$, continentur, de quarum tertia et quarta illud iam ostensum est), \emph{non-residuum vero omnium numerorum imparium qui ipsius 5 sint non-residua}, ut iam supra demonstravimus. Clarum autem est, hoc theorema sufficere ad diiudicandum, utrum $+5$ (eoque ipso, $-5$, si tamquam productum ex $+5$ et $-1$ consideretur) numeri cuiuscunque dati residuum sit an non-residuum. Denique observetur huius theorematis cum illo quod art.\,120 de residuo $-3$ exposuimus analogia.

At verificatio illius inductionis non adeo facilis. Quando numerus primus formae $20n+1$, sive generalius formae $5n+1$ proponitur, res simili modo absolvi potest, ut in artt.\,114, 119. Sit scilicet numerus quicunque pro modulo $5n+1$ ad exponentem 5 pertinens $a$, quales dari ex Sect.\,praec. manifestum, \SourceLeaf{97}eritque $a^5\equiv1$, sive $(a-1)(a^4+a^3+a^2+a+1)\equiv0\ (\mathrm{mod.}\,5n+1)$. At quia nequit esse $a\equiv1$, neque adeo $a-1\equiv0$; necessario erit $a^4+a^3+a^2+a+1$ $\equiv0$. Quare etiam $4(a^4+a^3+a^2+a+1)=(2aa+a+2)^2-5a^2$ erit $\equiv0$ \textit{i.\ e.} $5a^2$ erit residuum ipsius $5n+1$, adeoque etiam $5$, quia $a^2$ est residuum per $5n+1$ non divisibile ($a$ enim per $5n+1$ non divisibilis propter $a^5\equiv1)$. \textit{Q.\ E.\ D.}

At casus, ubi numerus primus formae $5n+4$ proponitur, subtiliora artificia postulat. Quoniam vero propositiones quarum ope negotium absolvitur in sequentibus generalius tractabuntur, hic breviter tantum eas attingimus.

I. Si $p$ est numerus primus atque $b$ non-residuum quadraticum datum ipsius $p$, valor expressionis

\[
(A)\ldots\quad\frac{(x+\surd b)^{p+1}-(x-\surd b)^{p+1}}{\surd b}
\]

(ex qua evoluta irrationalitatem abire facile perspicitur), semper per $p$ divisibilis erit, quicunque numerus pro $x$ assumatur. Patet enim ex inspectione coefficientium qui ex evolutione ipsius $A$ obtinentur, omnes terminos a secundo usque ad penultimum (incl.) per $p$ divisibiles fore, adeoque esse $A\equiv2(p+1)(x^p+xb^{\frac{p-1}{2}})$ $(\mathrm{mod.}\,p)$. At quoniam $b$ ipsius $p$ non-residuum est, erit $b^{\frac{p-1}{2}}\equiv-1\ (\mathrm{mod.}\,p)$, (art.\ $106$); $x^p$ autem semper est $\equiv x$ (Sect.\ praec.), unde fit $A\equiv0$. \textit{Q.\ E.\ D.}

II. In congruentia $A\equiv0\ (\mathrm{mod.}\,p)$, indeterminata $x$ habet $p$ dimensiones omnesque numeri $0,1,2....p-1$ illius radices erunt. Iam ponatur $e$ esse divisorem ipsius $p+1$ eritque expressio

\[
\frac{(x+\surd b)^e-(x-\surd b)^e}{\surd b}
\]

(quam per $B$ designamus) si evolvitur, ab irrationalitate libera, indeterminata $x$ in ipsa $e-1$ dimensiones habebit. constatque ex analyseos primis elementis, $A$ per $B$ (indefinite) esse divisibilem. Iam dico $e-1$ valores ipsius $x$ dari, quibus in $B$ substitutis, $B$ per $p$ divisibilis evadat. Ponatur enim $A\equiv BC$, habebitque $x$ in $C$ dimensiones $p-e+1$, adeoque congruentia $C\equiv0\ (\mathrm{mod.}\,p)$ non plures quam $p-e+1$ radices. Unde facile patet, omnes reliquos numeros ex his $0,1,2,3....p-1$, quorum multitudo $=e-1$, congruentiae $B\equiv0$ radices fore.

III. Iam ponatur $p$ esse formae $5n+4$, $e=5$, $b$ non-residuum ipsius $p$, atque numerum $a$ ita determinatum. ut sit

\SourceLeaf{98}\[
\frac{(a+\surd b)^5-(a-\surd b)^5}{\surd b}
\]

per $p$ divisibilis. At illa expressio fit

\[
=10a^4+20aab+2bb=2((b+5aa)^2-20a^4)
\]

Erit igitur etiam $(b+5aa)^2-20a^4$ per $p$ divisibilis \textit{i.\ e.} $20a^4$ residuum ipsius $p$; at quoniam $4a^4$ residuum est per $p$ non divisibile (facile enim intelligitur, $a$ per $p$ dividi non posse), etiam $5$ residuum ipsius $p$ erit. \textit{Q.\ E.\ D.}

Hinc patet theorema in initio huius articuli prolatum generaliter verum esse. ---

Observamus adhuc, demonstrationes pro utroque casu ill.\ La Grange deberi, \textit{M\'em. de l'Ac. de Berlin} $1775$, \textit{p.}\ $352$ sqq.

\begin{center}\textit{De $\pm7$.}\end{center}

\Article{124}

Per similem methodum demonstratur,

$-7$ \textit{esse non-residuum cuiusvis numeri qui ipsius $7$ sit non-residuum.}

Ex inductione vero concludi potest,

$-7$ \textit{esse residuum cuiusvis numeri primi qui ipsius $7$ sit residuum.}

At hoc a nemine hactenus rigorose demonstratum. Pro iis quidem residuis ipsius $7$, quae sunt formae $4n-1$, facilis est demonstratio; etenim per methodum ex praecc.\ abunde notam ostendi potest, $+7$ semper esse talium numerorum primorum non-residuum, adeoque $-7$ residuum. Sed parum hinc lucramur: reliqui enim casus per hanc methodum tractari nequeunt. Unum quidem adhuc casum simili modo ut artt.\ $119$, $123$ absolvere possumus. Scilicet si $p$ est numerus primus formae $7n+1$, atque $a$ pro modulo $p$ ad exponentem $7$ pertinens. facile perspicitur

\[
\frac{4(a^7-1)}{a-1}=(2a^3+a^2-a-2)^2+7(a^2+a)^2
\]

per $p$ divisibilem, adeoque $-7(a^2+a)^2$ ipsius $p$ residuum fore. At $(a^2+a)^2$. tamquam quadratum, ipsius $p$ residuum est, insuperque per $p$ non divisibile; quum enim $a$ ad exponentem $7$ pertinere supponatur, neque $\equiv0$, neque $\equiv-1\ (\mathrm{mod.}\,p)$ esse potest, \textit{i.\ e.} neque $a$ neque $a+1$ per $p$ divisibilis erit, adeoque etiam quadratum $(a+1)^2a^2$. Unde manifesto etiam $7$ ipsius $p$ residuum \SourceLeaf{99}erit. \textit{Q.\ E.\ D} --- At primi numeri formae $7n+2$ vel $7n+4$ omnes methodos hucusque traditas eludunt. Ceterum etiam haec demonstratio ab ill.\ La Grange primum est detecta l.\ c.\ --- Infra Sect.\ VII.\ docebimus generaliter, expressionem $\frac{4(x^p-1)}{x-1}$ semper ad formam $X^2\mp pY^2$ reduci posse, (ubi signum superius est accipiendum quando $p$ est numerus primus formae $4n+1$, inferius quando est formae $4n+3$), denotantibus $X,Y$ functiones rationales ipsius $x$, a fractionibus liberas. Hanc discerptionem ill.\ La Grange ultra casum $p=7$ non perfecit v.\ l.\ c.\ p.\ $352$.

\begin{center}\textit{Praeparatio ad disquisitionem generalem.}\end{center}

\Article{125}

Quoniam igitur methodi praecedentes ad demonstrationes generales stabiliendas non sufficiunt, iam tempus est, aliam ab hoc defectu liberam exponere. Initium facimus a theoremate, cuius demonstratio satis diu operam nostram elusit, quamvis primo aspectu tam obvium videatur, ut quidam ne necessitatem quidem demonstrationis intellexerint. Est vero hoc: \textit{Quemvis numerum, praeter quadrata positive sumta, aliquorum numerorum primorum non-residuum esse.} Quia vero hoc theoremate tantummodo tamquam auxiliari ad alia demonstranda usuri sumus, alios casus hic non explicamus quam quibus ad hunc finem indigemus. De reliquis casibus postea sponte idem constabit. Ostendemus itaque, \textit{quemvis numerum primum formae $4n+1$, sive positive sive negative accipiatur}\HistoricalNote{*)}{$+1$ autem excipi oportere per se manifestum est.}, \textit{non-residuum esse aliquorum numerorum primorum}, et (si $>5$) quidem talium qui ipso sint minores.

\textit{Primo}, quando numerus primus $p$, formae $4n+1$ ($>17$; sed $-13\,N3$, $-17\,N5$), \textit{negative} sumendus proponitur, sit $2a$ numerus par proxime maior quam $\surd p$; tum facile perspicitur, $4aa$ semper fore $<2p$ sive $4aa-p<p$. At $4aa-p$ est formae $4n+3$, $+p$ autem residuum quadraticum ipsius $4aa-p$, (quoniam $p\equiv4aa\ (\mathrm{mod.}\,4aa-p)$); quodsi igitur $4aa-p$ est numerus primus, $-p$ ipsius non-residuum erit; sin minus, necessario factor aliquis ipsius $4aa-p$ formae $4n+3$ erit; et quum $+p$ etiam huius residuum esse debeat, $-p$ ipsius non-residuum erit. \textit{Q.\ E.\ D.}

Pro numeris primis \textit{positive} sumendis duos casus distinguimus. \textit{Primo} sit $p$ numerus primus formae $8n+5$. Sit $a$ numerus quicunque positivus $<\surd\!\frac12 p$. Tum $8n+5-2aa$ erit numerus positivus formae $8n+5$ vel $8n+3$ (prout $a$ \SourceLeaf{100}par vel impar), adeoque necessario per numerum aliquem primum formae $8n+3$ vel $8n+5$ divisibilis, productum enim ex quotcunque numeris formae $8n+1$ et $8n+7$ neque formam $8n+3$ neque hanc $8n+5$ habere potest. Sit hic $q$, eritque $8n+5\equiv2a^2\ (\mathrm{mod.}\,q)$. At $2$ ipsius $q$ non-residuum erit (art. $112$), adeoque etiam $2a^2$\HistoricalNote{*)}{Art. $98$. Patet enim $a^2$ esse residuum ipsius $q$ per $q$ non divisibile, nam alias etiam numerus primus $p$ per $q$ foret divisibilis. \textit{Q. E. A.}} et $8n+5$. \textit{Q.\ E.\ D.}

\Article{126}

Sed numerum quemvis primum formae $8n+1$ positive acceptum semper alicuius numeri primi ipso minoris non residuum esse, per artificia tam obvia demonstrari nequit. Quum autem haec veritas maximi sit momenti, demonstrationem rigorosam, quamvis aliquantum prolixa sit, praeterire non possumus. Praemittimus sequens

\textsc{Lemma.} \textit{Si habentur duae series numerorum}

\[
A,B,C\ \textit{etc.}\ ....\ (\mathrm{I}),\quad A^{\prime},B^{\prime},C^{\prime}\ \textit{etc.}\ ....\ (\mathrm{II})
\]

(utrum terminorum multitudo in utraque eadem sit necne nihil interest) \textit{ita comparatae, ut, denotante $p$ numerum quemcunque primum aut numeri primi potestatem, terminum aliquem secundae seriei (sive etiam plures) metientem, totidem ad minimum termini in serie prima sint per $p$ divisibiles, quot sunt in secunda: tum dico productum ex omnibus numeris} (I) \textit{divisibile fore per productum ex omnibus numeris} (II).

\textit{Exempl.} Constet (I) e numeris $12,18,45$; (II) ex his $3,4,5,6,9$. Tum divisibiles erunt per $2,4,3,9,5$ in (I) $2,1,3,2,1$ termini, in (II) $2,1,3,1,1$ termini, respective; productum autem omnium terminorum (I) $=9720$ divisibile est per productum omnium terminorum (II), $3240$.

\textit{Demonstr.} Sit productum ex omnibus terminis (I), $=Q$, productum omnium terminorum seriei (II), $=Q'$. Patet quemvis numerum primum qui sit divisor ipsius $Q'$ etiam ipsius $Q$ divisorem fore. Iam ostendemus quemvis factorem primum ipsius $Q'$, in $Q$ totidem ad minimum dimensiones habere quot habeat in $Q'$. Esto talis divisor $p$, ponaturque, in serie (I) $a$ terminos esse per $p$ divisibiles. $b$ terminos per $p^2$ divisibiles, $c$ terminos per $p^3$ divisibiles etc., similia denotent literae $a',b',c'$ etc. pro serie (II), perspicieturque facile, $p$ in $Q$ habere \SourceLeaf{101}$a+b+c+$ etc. dimensiones, in $Q'$ vero $a'+b'+c'+$ etc. At $a'$ certe non maior quam $a$. $b'$ non maior quam $b$ etc. (hyp.); quare $a'+b'+c'$ etc. certo non erit $>a+b+c$ etc. --- Quum itaque nullus numerus primus in $Q'$ plures dimensiones habere possit, quam in $Q$, $Q$ per $Q'$ divisibilis erit (art. $17$). \textit{Q.\ E.\ D.}

\Article{127}

\textsc{Lemma.} \textit{In progressione $1,2,3,4,....n$, plures termini esse nequeunt per numerum quemcunque $h$ divisibiles, quam in hac $a,a+1,a+2.....a+n-1$ ex totidem terminis constante.}

Nullo enim negotio perspicitur si $n$ fuerit multiplum ipsius $h$, in utraque progressione $\frac nh$ terminos fore per $h$ divisibiles; sin minus, ponatur $n=eh+f$, ita ut $f$ sit $<h$, eruntque in priori serie $e$ termini per $h$ divisibiles, in posteriori autem vel totidem vel $e+1$.

Hinc tamquam Coroll. sequitur propositio ex numerorum figuratorum theoria nota, sed a nemine, ni fallimur, hactenus directe demonstrata,

\[
\frac{a.a+1.a+2....a+n-1}{1\ .\ 2\ .\ 3\ ....\ n}
\]

semper esse numerum integrum.

Denique Lemma hoc generalius ita proponi potuisset:

In progressione $a,a+1,a+2,....a+n-1$ totidem ad minimum dantur termini secundum modulum $h$ numero cuicunque dato, $r$, congrui, quot in hac $1,2,3....n$ termini per $h$ divisibiles.

\Article{128}

\textsc{Theorema.} \textit{Sit $a$ numerus quicunque formae $8n+1$, $p$ numerus quicunque ad $a$ primus, cuius residuum $+a$, tandem $m$ numerus arbitrarius: tum dico, in progressione}

\[
\begin{gathered}a,\ \tfrac12(a-1),\ 2(a-4),\ \tfrac12(a-9),\ 2(a-16),\\ ....\ 2(a-m^2),\quad\textit{vel}\quad\tfrac12(a-m^2)\end{gathered}
\]

\textit{prout $m$ par vel impar, totidem ad minimum dari terminos per $p$ divisibiles quot dentur in hac}

\[
1,\ 2,\ 3,\ ....\ 2m+1
\]

Priorem progressionem designamus per (I), posteriorem per (II).
